\documentclass[../main.tex]{subfiles}

\begin{document}

\ifSubfilesClassLoaded{

    \tableofcontents
    
    %\setcounter{chapter}{}
}{}

\chapter{ Introduction }
% 代靖涵 25120222201319

\section{Function Spaces and Quotient Spaces}

\begin{definition}{}{}
    Given two topological spaces $X, Y$ we denote by $Y^X$ the set of all continuous functions (maps) from $X$ to $Y$. This set is topologised by the so-called \dfntxt{compact-open-topology}, viz., the topology with the subbase as the collection of all sets of the form
\[
\langle K, U \rangle = \{f \in Y^X : f(K) \subset U\}
\]
where $K \subset X$ is compact and $U \subset Y$ is open.
\end{definition}

\begin{theorem}{Exponential correspondence}{10-4-2}
Let $X$ be a locally compact Hausdorff space and $Y,Z$ be any two topological spaces.
\begin{enumerate}
\item[(a)] The \dfntxt{evaluation map} $E:Y^X \times X \to Y$ given by $E(f,x) = f(x)$ is continuous.
\item[(b)] A function $g : Z \to Y^X$ is continuous iff the composite $E \circ (g \times \mathrm{Id}_X) : Z \times X \to Y$ is continuous.
\item[(c)] If $Z$ is also locally compact and Hausdorff, then the function
\[
\psi : (Y^X)^Z \to Y^{Z \times X}
\]
defined by
\[
\psi(g) = E \circ (g \times \mathrm{Id}_X)
\]
is a homeomorphism.
\end{enumerate}
\end{theorem}

\begin{note}
    \begin{itemize}
        \item (a) 是说, 我们可以把 \(  X\to Y  \)的连续性像讲废话一样, 特殊地多复制一份放到左边, 变成 \(  Y^{X}  \times X\to Y\)的连续性.  也就是说如果\(  X\to Y  \)是连续的, 我们可以预先准备一个 \(  X  \)中的一个东西, 然后让 \(  Y^{X}  \)中的家伙把 \(  X \)   里面的东西给吃了, 这个操作也是连续的.
        \item (b) 我们发现抽象的紧开拓扑的连续性可以和具体乘积拓扑的连续性密切联系. 形象地说, 我们可以把右边的指数拽成左边的乘积. 虽然 \(  Z\times X\to Y  \)上的连续性, 体现出来地更破碎一点(关于每个 \(  \left( x_0,k \right)   \)附近的``块状邻域'', 而不是 \(  Z\to Y^{X}  \)连续所需要的 像是包裹了 \(  \left\{ x_0 \right\}\times K  \)一样的长条邻域. ). 但是由于 \(  K  \)的紧性, 我们可以将``块状邻域''拼成``长条状''的. 
    \end{itemize}
    
\end{note}
\begin{proofsketch}
    (c) :证明映射连续, 可以利用 (b)把指数上的东西不断地拽到左边来, 直到右边没有指数了. 发现得到的东西恰好是``依次预先准备左边要的贡品, 然后让一个家伙依次把这些贡品吃进去, 直到没有东西可以吃了'', 这个操作根据我们(a)中的note是连续的.
\end{proofsketch}
\begin{proof}
    (a) Take any open subset \(  U  \) of \(  Y  \). We need to show that \(  E^{-1} \left( U \right)   \) is open. For which, we take \(  \left( f, x_0 \right)\in E^{-1} \left( U \right)   \). Since \(  X  \) is locally compact, there exists a compact neighbourhood \(  K  \)\footnote{That is , a compact set \(  K  \) such that there exists open set \(  V  \) , \(  x\in V\subseteq K  \)   } of \(  x_0  \) contained in \(  f^{-1} \left( U \right)   \) , that is \(  f\in \left<K,U \right>  \) .   Then for each \( \left( g,x_1 \right) \in \left<K,U \right>\times K   \), we have \[
    E\left(  g,x_1\right)= g\left( x_1 \right)\in U  
    \] which implies that  \(  \left( f,x_0 \right)\in \left<K,U \right>\times K\subseteq E^{-1} \left( U \right)     \). Thus \(  E  \) is continuous.   
    
    (b) From (a) we only need to show that if \(  E\circ \left( g\times \mathrm{Id}_{X} \right)   \) is continuous, then so does \(  g  \). For which take any subbasis \(  \left<K,U \right>   \) of \(  Y^{X}  \), we only need to show that \(  g^{-1} \left<K,U \right>  \) is open. Take any \(  z_0 \in g^{-1} \left( K,U \right)   \) ,      then \(  g\left( z_0 \right) \in \left<K,U \right>   \). For each \(  k \in K  \), we have \[
    E\circ \left( g \times \mathrm{Id}_{X} \right)\left(z_0,k\right)  = g\left( z_0 \right)\left( k \right)\in U  
    \]  Since \(  E\circ \left( g\times \mathrm{Id}_{X} \right)   \)  is continuous, there exists open neighbourhood \(  V_{k}  \) of \(  z_0  \)   and the open neighbourhood \(  W_{k}  \) of \(  k  \) such that  \[
    E\circ \left( g\times \mathrm{Id}_{X} \right)\left( V_{k}\times W_{k} \right)\subseteq U  
    \]  Since \(  K  \) is compact, we can pass on to a finite open cover \(  K= \bigcup _{i= 1}^{m}W_{k_{i}} \) , and set \(  V\coloneqq \bigcap _{i= 1}^{m}V_{k_{i}}  \). Then \[
    g\left( V \right)\left( K \right)= E\circ \left( g\times \mathrm{Id}_{X} \right)\left( V\times K \right)\subseteq U\implies g\left( V \right)\subseteq \left<K,U \right>     
    \]   Thus \(  g  \) is continuous. 

    (c) Note that from (b) it follows that \(  \psi   \) is  well-defined and a bijection. Applying (b) with \(  \left( Y^{X} \right)^{Z}   \)in place of \(  Z  \)  and  \(  Y^{Z\times X}  \) place of \(  Y^{X}  \). Then the continuity of  \(  \psi   \) is equivalent to that of \(  E\circ \left( \psi \times \mathrm{Id}_{Z\times Y} \right):\left( Y^{X} \right)^{Z}\times \left( Z\times X \right) \to Y    \)     given by \[
    \left(  \lambda ,z,x \right)\mapsto  \lambda \left( z \right)\left( x \right)      \] This latter map is the composite of the two maps \[
    \left( \left( Y^{X} \right)^{Z}\times Z  \right)\times X\to Y^{X}\times X\to Y 
    \]The continuity of this two maps follows from (a).

    The continuity of \(  \psi   ^{-1} \) is equivalent to that of  \(  E\circ \left( \psi ^{-1} \times \mathrm{Id}_{Z} \right)   : \left( Y^{Z\times X}\times Z \right)\to Y^{X} \) , and is quivalent to that of \(  E\circ \left( E\circ \left( \psi ^{-1} \times \mathrm{Id}_{Z} \right)\times \mathrm{Id}_{X}  \right):\left( Y^{Z\times X}\times Z \right)\times X\to Y    \), given by \[
    \left( \mu ,z,x \right)\mapsto \mu \left( z,x \right)  
    \] This latter map is the evaluation map\[
   Y^{Z\times X}\times \left( Z\times X \right)\to Y  
    \] \[
    E\left( \mu ,\left( z,x \right)  \right)= \mu \left( z,x \right)  
    \] which is continuous from (a).
\end{proof}

\subsection{Quotient Space}
\begin{lemma}{}{}
The following statements are equivalent.
\begin{enumerate}
    \item[(i)] $U \in \tau'$ iff $q^{-1}(U) \in \tau$.
    \item[(ii)] A function $g : (Y, \tau') \to (Z, \tau'')$ is continuous iff $g \circ q$ is continuous.
    \item[(iii)] For a fixed $\tau$, $\tau'$ is the maximal topology on $Y$ such that $q$ is continuous.
\end{enumerate}
\end{lemma}

\begin{definition}{}{}
Under the above conditions, we say $(Y, \tau')$ is a \dfntxt{quotient space} of $(X, \tau)$ and the map $q$ is called a \dfntxt{quotient map}.
\end{definition}

\begin{proposition}{}{}
    \(  Y  \) is a \(  T_1  \) space iff all the fibers of \(  q  \) are closed subsets, irrespective of whether \(  X  \) is \(  T_1  \) or not.     
\end{proposition}
\begin{proposition}{}{}
    If \(  q:X\to Y  \) is a surjective open mapping given by a closed relation \(  R\subseteq X\times X  \). Then \(  Y  \) is a Hausdorff space, irrepective of whether \(  X  \) is Hausdorff or not.    
\end{proposition}
\begin{proof}
    \(  Y  \) is Hausdorff, iff \(   \Delta _{Y}  \) is closed. We know from definition that  \[
    \left( q \times q \right)^{-1} \left(  \Delta _{Y} \right)= R  
    \]  Then \[
    R^{c}= \left( q \times  q \right)^{-1} \left(   \Delta _{Y}^{c} \right)\implies \left( q\times q \right)\left( R^{c} \right)=  \Delta _{Y}^{c}     
    \]Since \(  q  \) is a open map, so dose \(  q\times q  \), and since \(  R^{c}  \) is open, so dose \(   \Delta _{Y}^{c}  \).    Thus \(   \delta _{Y}  \) is closed,  \(  Y  \) is Hausdorff.  
\end{proof}

\begin{theorem}{}{}
Let $q : X \to Y$ be a quotient map and $A \subset X$. Then the restriction map $q_A : A \to q(A) =: B$ is a quotient map if the following conditions are satisfied:
\begin{enumerate}
\item[(a)] $B$ is the intersection of an open set and a closed set in $Y$.
\item[(b)] For $C \subset B$, $q_A^{-1}(C)$ is closed in $A$ (respectively open in $A$) iff $q^{-1}(C)$ is closed (respectively open) in $q^{-1}(B)$.
\end{enumerate}
\end{theorem}

\begin{proof}
    It is obvious that \(  q  \) is continuous and surjective. 
   Let \(  C  \) be a closed sets in \(  B  \). We need to prove that \(  q_{A}^{-1} \left( C \right)   \) is closed in \(  A  \)    . We write \(  C= K\cap B  \), where \(  K  \) is a closed subset of \(  X  \).   \(  C  \) is closed in \(  X  \). Then \(  q^{-1} \left( C \right)    \) is closed in \(   X  \).    We know from \[
   q^{-1} \left( C \right)=  q^{-1} \left( C \right)\cap q^{-1} \left( B \right)   
   \] that \(  q^{-1} \left( C \right)    \) is closed in \(  q^{-1} \left( B \right)   \).   Thus it follows from (b) that \(  q_{A}^{-1} \left( C \right)   \) is closed in \(  A  \).   
   
    Then we suppose that \(  q_{A}^{-1} \left( C \right)   \) is closed in \(  A  \), then \(  q^{-1} \left( C \right)    \) is closed in \(  q^{-1} \left( B \right)   \).    We know from (a) that there is a open set \(  U  \) and a closed set \(  F  \) such that \[
   B= U\cap F
   \]  Then we have \[
   q^{-1} \left( B \right)=  q^{-1} \left( U \right)\cap  q ^{-1} \left( F \right)   
   \]Where \(  q^{-1} \left( U \right)   \) is open, \(  q^{-1} \left( F \right)   \) is closed.   There is a closed set \(  K^{\prime}   \) of \(  X  \) such that \[
   q^{-1} \left( C \right)= K^{\prime} \cap q^{-1} \left( B \right)=  q^{-1} \left( U \right)\cap \left( q^{-1} \left( F \right)\cap K^{\prime}   \right)    
   \]  Let \(  \tilde{K}=  q^{-1} \left( F \right)\cap K^{\prime}    \), then \(  \tilde{K}  \) is closed in \(  X  \).   Then \[
   q^{-1} \left( U\setminus C \right)=  q ^{-1} \left( U \right)\setminus   q ^{-1} \left( C \right)=  q ^{-1} \left( U \right)\setminus \tilde{K}^{\prime}     
   \] which is open in \(  X  \). Thus \(  U\setminus C  \) is open in \(  Y  \).    \[
   \left( U\setminus C \right)^{c}\cap B= \left( U\cap C^{c} \right)^{c}\cap B= \left( U^{c}\cup C \right) \cap B = C\cap B= C
   \]Where \(  \left( U\setminus C \right)^{c}   \) is closed in \(  Y  \). Thus \(  C  \) is closed in \(  B  \).
   
   The remaining part is similarly.

\end{proof}
   \begin{theorem}{}{}
If $Y$ is a locally compact Hausdorff space, then for any quotient map $q:X \to Z$, the product $q \times \mathrm{Id}_Y$ is a quotient map.
\end{theorem}
\begin{proofsketch}
    为了利用原本的商映射的泛性质, 我们需要把冗余的乘积 \(  \times Y  \)丢到右边的指数上去, 这个操作就需要我们甩过去的 \(  Y  \)是局部紧的Hausdorff空间.  
\end{proofsketch}
\begin{proof}
    Take any map \(  g:  Z\times Y\to W  \). Let \(  f\coloneqq g\circ \left( q\times \mathrm{Id}_{Y} \right)  :X\times Y\to W \).   If \(  f  \) is continuous, define \[
    \hat{f}: X\to W^{Y},
    \] \[
    \hat{f}\left( x \right)\left( y \right)= f\left( x,y \right)
    \]  Then \[
    E\circ \left( \hat{f}\times \mathrm{Id}_{Y} \right): X\times Y\to W 
    \] \[
    E\circ \left( \hat{f}\times \mathrm{Id}_{Y} \right)\left( x,y \right)= \hat{f}\left( x \right)\left( y \right)  = f\left( x,y \right)   
    \]Then it follows from Theorem \ref{thm:10-4-2} that \(  \hat{f}  \) is continuous. There exists a continuous function \(  \tilde{f}: Z\to W^{Y}  \), such that \(  \tilde{f}\circ q= \hat{f}  \).  Then \(  E\circ \left( \tilde{f} \times \mathrm{Id}_{Y} \right)   \) is continuous, we have for \(  z \in Z  \), take \(  x \in X  \) such that \(  z=  q\left( x \right)   \), then    \[
   \begin{aligned}
    E\circ \left( \tilde{f}\times \mathrm{Id}_{Y} \right)\left( z,y \right)= \tilde{f}  \left( z \right)\left( y \right)  = \tilde{f}\left( q\left( x \right)  \right)\left( y \right)= \hat{f}\left( x \right)\left( y \right)= f\left( x,y \right)    
   \end{aligned}
    \] The other side,  \[
    g\left( z,y \right)= g\left( q\left( x \right),y  \right)  = g\circ \left( q\times \mathrm{Id}_{Y} \right)\left(x,y \right)= f\left( x,y \right)   
    \] We have \[
    g = E\circ \left( \tilde{f}\times \mathrm{Id}_{Y} \right) 
    \] is continuous.
\end{proof}

\begin{corollary}{}{}
If $q_i: X_i \to Z_i, i=1,2$ are quotient maps such that $X_1, Z_2$ or $Z_1, X_2$ are locally compact Hausdorff spaces, then $q_1 \times q_2$ is a quotient map.
\end{corollary}
\begin{proof}
    If \(  X_1,Z_2  \) are locally compact Hausdorff spaces. Then \[
    q_1\times q_2= \left( q_1\times \mathrm{Id}_{Z_2} \right)\circ \left( \mathrm{Id}_{X_1}\times q_2 \right) :X_1\times X_2\to Z_1\times Z_2
    \] is a quotient map.
\end{proof}

\section{Relative Homotopy}

\begin{definition}{}{}
Let $A \subset X$, and $f,g : X \longrightarrow Y$ be any two maps such that $f(a) = g(a), \forall a \in A$. We say $f$ is homotopic to $g$ \dfntxt{relative to} $A$ if there is a homotopy $H$ from $f$ to $g$ such that $H(a,t) = f(a), \forall a \in A$. We write this
\[
f \simeq g, \text{ rel } A.
\]
\end{definition}
\begin{definition}{}{} 
Two topological pairs $(X, A), (Y, A)$ are said to have the same homotopy type if there exist maps $f : (X, A) \to (Y, A), g : (Y, A) \to (X, A)$ such that $f_A = Id_A = g_A$ and
\[
g \circ f \simeq Id_X \text{ rel } A; \quad f \circ g \simeq Id_Y \text{ rel } A.
\]
\end{definition}

\begin{definition}{}{}
    Let \(  X  \) be a topological space and \(  A\subseteq X  \) a subspace, \(  \iota :A\hookrightarrow X  \) the inclusion map. 
    
    \begin{itemize}
        \item By a \dfntxt{retraction} \(  r:X\to A  \), we mean a map \(  r  \) such that \(  r\left( a \right)= a, \forall a \in  A  \), equivalently, \(  r\circ \iota = \operatorname{Id}_{A}  \).    
        \item A map \(  f:X\to X  \) which is homotopic to \(  \operatorname{Id}_{X}  \) is called a \dfntxt{deformation }.
        \item If \(  r :X\to A \) is a rectraction   , \( \iota \circ r  \) is homotopoc to \(  \operatorname{Id}_{X}  \), then we call \(  r  \) a \dfntxt{deformation rectraction}, and \(  A  \) a \dfntxt{deformation recract(DR)} of \(  X  \) .
        \item If \(  r:X\to A  \) is a \dfntxt{deformation rectraction}, and the homotopy from \(  \operatorname{Id}_{X}  \) to \(  r  \) is relative to \(  A  \), then \(  r  \) is called a \dfntxt{strong deformation rectraction(SDR)}.     
        \item If \(  \iota :A\hookrightarrow X  \) is a homotopy equivalence , then we say \(  A  \) is a \dfntxt{week deformation rectract (WDR)} of \(  X  \)   
    \end{itemize}
    
\end{definition}

\begin{remark}
    SDR\(  \implies   \) DR \(  \implies   \) WDR  
\end{remark}

\begin{example}{}{}
    \begin{enumerate}
        \item Let \(  X  \) be a convex space, then every point \(  x_0 \in X  \) is a \(  SDR  \) of \(  X  \).    We define a map \(  H: X\times \mathbb{I}\to X  \),  \[
        H\left( x,t \right)= \left( 1-t \right)x+ tx_0  
        \]which is well defined since \(  X  \) is convex. Then \[
         H: \operatorname{Id}_{X}\simeq  c_{x_0}   , \operatorname{rel} \left\{ x_0 \right\}
        \]
    \end{enumerate}
    
\end{example}

\section{Typical Constructions}
\subsection{The Cone}

\begin{definition}{}{} 
By a cone $CX$ over $X$, we mean the quotient space of $X \times I$ obtained by identifying the subspace $X \times 0$ to a single point. (See Figure 1.18.) This point is called the \dfntxt{vertex} or \dfntxt{apex} of the cone and is denoted simply by $\star$. The space $X$ itself can then be identified with the image of the subspace $X \times 1$ in the quotient space, and this is then referred to as the \dfntxt{base} of the cone.
\end{definition}

\begin{remark}
    
\begin{itemize}
    \item One can represent of point of \(  CX  \) by \(  \left[ x,t \right]   \) for some \(  x \in X  \) and \(  t\in \mathbb{I}  \). The representation is unique  if \(  t= 0  \). Also, \(  \left[ x,0 \right]   \) represents the vertex \(  \star  \).       
    \item \(  \left( [x,t],s \right)\to \left[ x,st \right]    \) defines a homotopy of the constant map \(  \star  \) with the identify map of \(  CX  \), then \(  CX  \) is contractible.    
\end{itemize}

\end{remark}

\begin{definition}{}{}
    Given any map \(  f:X\to Y  \), there is a map \(  C\left( f \right):CX\to CY   \) defined by \[
     C\left( f \right)\left[ x,t \right]= \left[ f\left( x \right),t  \right]  
    \]satisfying \[
\begin{tikzcd}
    X \arrow[r, "f"] \arrow[d, "i_X"'] & Y \arrow[d, "i_Y"] \\
    CX \arrow[r, "C(f)"] & CY
\end{tikzcd}
\] where \(  i_{X}: X\to X\times \left\{ 1 \right\}  \) , \(  i_{Y}: Y\to Y\times \left\{ 1 \right\}  \) are the identifications. 
\end{definition}

\begin{remark}
    We always omitt the map \(  i_{X}  \) and simply write \(  X= X\times \left\{ 1 \right\}  \).  
\end{remark}

\begin{proposition}{}{}
    \(  C\left( f \right)   \) is a homeomorphism iff \(  f  \) is a homeomorphism.  
\end{proposition}
\begin{proof}
    If \(  C\left( f \right)   \) is a homeomorphism , then \(  C\left( f \right)|_{X}: X\to C\left( f \right)\left( X \right)     \)   is a homeomorphism with the subspace topology, where \[
    C\left( f \right)\left( X \right)= \left[ f\left( X \right),1  \right]= \left[ Y,1 \right]= Y    
    \]Thus \(  C\left( f \right)|_{X}:X\to Y   \) is a homeomorphism .

    Conversely , if \(  f  \) is a homeomorphism , then \(  f, f^{-1}   \) give rise to   continuous maps \(  C\left( f \right)   \) and \(  C\left( f^{-1}  \right)   \). Note that \[
    C\left( f \right)\circ C\left( f^{-1}  \right)\left( [y,t] \right)= [y,t],\quad C\left( f^{-1}  \right)\circ C\left( f \right)\left( [x,t] \right)= [x,t]      
    \]  We know \(  C\left( f \right)   \) is a homeomorphism. 
\end{proof}


\begin{theorem}{}{}
Any topological space $X$ is contractible iff it is a retract of the cone $CX$.
\end{theorem}

\begin{note}
  我们可以把 \(  X  \)在时间\(  \mathbb{I}  \)上的状态展开在    \(  X\times \mathbb{I}  \)上面. 
\end{note}

\begin{proof}
    Let \(  q:X\times \mathbb{I}\to CX  \) denote the quotient map. Given a homotopy \(  H\times \mathbb{I}\to X  \), such that \(  H\left( x,0 \right)= x_0   \), \(  H\left( x,1 \right)= x   \). There is a continuous map \(  r:CX\to X  \) such that \(  H= r\circ q  \). Clearly \(  r\left( x,1 \right)= x   \) and since we have identified \(  X  \) with \(  X\times 1  \), this means \(  X\times 1\subseteq CX  \) is a retract of \(  CX  \).
    
    Conversely, given any retraction \(  r:CX\to X  \) , the formula \(  H\left( x,t \right)= r\left( [x,t] \right)    \) gives a homotopy of the constatn map with the identity map of \(  X  \) and hence \(  X  \) is contractible.    
\end{proof}

\begin{theorem}{}{}
    Let \(  X,Y  \) be topological spaces. Then \(  f:X\to Y  \) is null   homotopic , if and only if \(  f  \) can be extend to  a map \(  \bar{f}:CX\to Y  \) 
\end{theorem}

\begin{proofsketch}
    \(  \bar{f}  \)就是同伦映射 \(  H: X\times \mathbb{I}\to Y  \)在商空间上的诱导. 
\end{proofsketch}

\begin{proof}
    If \(  f   \) is a null homotopic, then there exists \(  H:X\times \mathbb{I}\to Y  \) , \(  H\left( x,1 \right)= f\left( x \right)    \), \(  H\left( x,0 \right)= c   \). Then \(  H  \) induces a map on quotient \( \bar{f}:CX\to Y  \), such that \(  \bar{f}\left( x,1 \right)= f\left( x \right)    \). 

    Conversely, if \(  f  \) can be extend to \(  \bar{f}:CX\to Y  \). Then we define \[
    H\left( x,t \right)= \bar{f}\left( [x,t] \right)  
    \] We have \[
    H\left( x,0 \right)= \bar{f}\left( \star \right),\quad H\left( x,1 \right)= \bar{f}\left( [x,1] \right)= f\left( x \right)     
    \]  
\end{proof}

\begin{example}{}{}
    The cone over \(  \mathbb{S}^{n-1}  \) is homeomorphic to \(  \mathbb{D}^{n}  \)  . Consider the map \(  P:\mathbb{S}^{n-1}\times \mathbb{I}\to \mathbb{D}^{n}  \) given by \(  \left( x,t \right)\mapsto tx   \). The map is clearly a continuous bijection and one-to one except that all points of \(  \bar{P}:C\mathbb{S}^{n-1}\to \mathbb{D}^{n}  \).   Since the domain is compact and the range is Hausdorff, \(  \bar{P}  \) is a homeomorphism.

\end{example}

\begin{theorem}{}{12-2-1}
Let $X$ be any topological space and $f\colon S^{n-1} \to X$ be any map. The following statements are equivalent:
\begin{enumerate}
\item $f\colon S^{n-1} \to X$ is null homotopic;
\item $f$ can be extended to a map $\bar f\colon \mathbb{D}^n \to X$;
\item $f$ is null homotopic relative to any given point $p \in S^{n-1}$.
\end{enumerate}
\end{theorem}
\begin{remark}
    这个定理就是在说, 定端零伦就是自由零伦.
\end{remark}
\begin{proof}
    The equivalence of 1.2. is immediate.

    \(  2.\implies 3.  \) Fix \(  p \in S^{n-1}  \subseteq \mathbb{D}^{n}\). Let \(  F: \mathbb{D}^{n}\times  \mathbb{I}\to \mathbb{D}^{n}  \) , \[
    F\left( x,t \right)= \left( 1-t \right)x+ tp  
    \] Then \(  \bar{f}\circ F|_{{S}^{n-1}\times \mathbb{I}} :S^{n-1}\times \mathbb{I}\to X  \)  is a homotopy relative to \(  p  \) from \(  f  \) to \(  c_{p}  \).
    
    \(  3.\implies  1. \) is obvious. 
\end{proof}


\subsection{The Ajunction Spaces}

\begin{definition}{Adjunction space}{}
Let $X$ be a closed subspace of a space $Z$ and $f:X\to Y$ be a continuous map. The \dfntxt{adjunction space} $Z\cup_{f}Y=:A_{f}$ of $f$ is defined to be the quotient of the disjoint union $Y\sqcup Z$ by the identification $x\sim f(x)$ for all $x\in X$. This is also called the \dfntxt{space obtained by attaching $Z$ to $Y$ via the map $f$}. Note that $Y$ can be identified with a closed subspace of $A_{f}$ via $j:y\mapsto [y]$. If the image of $f$ is closed in $Y$, then the image of $Z$ is closed in $A_{f}$. Also, if $f$ is injective then $Z$ can be identified with its image in $A_{f}$.
\end{definition}

\begin{proof}
  Let \(  q: Y\sqcup Z\to A_{f}  \) be the quotient map. Then \[
  q^{-1} \left( j\left( Y \right)  \right)=Y \sqcup X  
  \] 
\end{proof}
\end{document}